\documentclass[11pt,a4paper]{article}
\usepackage{amsmath,amssymb}
\usepackage{listings}
\usepackage{hyperref}
\usepackage[margin=1in]{geometry}
\usepackage{fancyhdr}

\pagestyle{fancy}
\fancyhf{}
\fancyhead[L]{Module 05: Arrays and Strings}
\fancyhead[R]{C Programming Course}
\fancyfoot[C]{\thepage}

\lstset{
  language=C,
  basicstyle=\ttfamily\small,
  keywordstyle=\bfseries,
  commentstyle=\itshape,
  numbers=left,
  numberstyle=\tiny,
  frame=single,
  breaklines=true,
  tabsize=4
}

\title{Module 05: Arrays and Strings}
\author{C Programming Course}
\date{}

\begin{document}
\maketitle
\thispagestyle{fancy}

\section{One-Dimensional Arrays}
\begin{lstlisting}
#include <stdio.h>

int main(void) {
    int nums[5] = {10, 20, 30, 40, 50};
    for (int i = 0; i < 5; i++) {
        printf("nums[%d] = %d\n", i, nums[i]);
    }
    return 0;
}
\end{lstlisting}

Arrays are zero-indexed. The size is fixed at compile time for stack-allocated arrays.

\section{Multi-Dimensional Arrays}
\begin{lstlisting}
#include <stdio.h>

int main(void) {
    int matrix[2][3] = {
        {1, 2, 3},
        {4, 5, 6}
    };
    for (int i = 0; i < 2; i++) {
        for (int j = 0; j < 3; j++) {
            printf("%d ", matrix[i][j]);
        }
        printf("\n");
    }
    return 0;
}
\end{lstlisting}

\section{Character Arrays and Strings}
In C, a string is a null-terminated character array.

\begin{lstlisting}
#include <stdio.h>

int main(void) {
    char greeting[] = "Hello";
    printf("%s has %d visible chars\n",
           greeting, 5);
    /* greeting[5] == '\0' */
    return 0;
}
\end{lstlisting}

\section{String Manipulation Without string.h}
\begin{lstlisting}
#include <stdio.h>

int my_strlen(const char *s) {
    int len = 0;
    while (s[len] != '\0') {
        len++;
    }
    return len;
}

int main(void) {
    char msg[] = "C Programming";
    printf("Length: %d\n", my_strlen(msg));
    return 0;
}
\end{lstlisting}

\section{Common String Operations}
\begin{itemize}
  \item \texttt{strlen} -- returns the length of a string.
  \item \texttt{strcpy} -- copies one string to another.
  \item \texttt{strcat} -- concatenates two strings.
  \item \texttt{strcmp} -- compares two strings lexicographically.
\end{itemize}

Always ensure destination buffers are large enough to prevent buffer overflows.

\end{document}
